\section*{Conclusiones}
\addcontentsline{toc}{section}{Conclusiones}

El desarrollo del plan de negocios para la plataforma EduFinanzas demuestra la viabilidad técnica y comercial de una solución orientada a cerrar la brecha de alfabetización financiera en Bogotá. A través de un diagnóstico multidimensional, se estableció que la problemática no radica únicamente en el acceso al sistema bancario, sino en la carencia de herramientas tecnológicas que permitan una gestión de recursos eficiente y personalizada para sectores tradicionalmente excluidos.

Desde la perspectiva estratégica, la implementación del modelo de negocio CANVAS permitió estructurar una propuesta de valor diferenciada, donde la tecnología actúa como el eje integrador entre la educación y la operatividad financiera. La identificación de segmentos críticos, como trabajadores informales y nómadas digitales, facilitó el diseño de una arquitectura de servicios capaz de adaptarse a flujos de ingresos variables, garantizando que la solución sea pertinente para la realidad socioeconómica del entorno colombiano.

%Verificar este parrafo con la seccion de analisis financiero%
El análisis financiero proyectado confirma la sostenibilidad del emprendimiento a largo plazo. Si bien la fase inicial presenta un flujo de caja negativo debido a la inversión en infraestructura tecnológica y adquisición de usuarios, el modelo de escalabilidad horizontal permite alcanzar el punto de equilibrio en un periodo razonable. La naturaleza exponencial de las proyecciones se sustenta en el bajo costo marginal de distribución de los servicios digitales y en una arquitectura diseñada para evolucionar ágilmente según las demandas del mercado.

Los instrumentos de diagnóstico estratégico, específicamente las matrices DOFA y PEYEA, sitúan al proyecto en una postura agresiva, propicia para la innovación en un mercado Fintech altamente competitivo. Se concluye que la principal fortaleza del proyecto reside en su sólido componente técnico, basado en microservicios y la capacidad de pivotaje metodológico de una startup, lo cual mitiga los riesgos asociados a la obsolescencia tecnológica y permite una respuesta inmediata ante los cambios en el modelo de negocio o en las sugerencias de los usuarios.

Finalmente, la construcción del prototipo mínimo viable valida la factibilidad técnica de los requerimientos identificados. El uso de patrones arquitectónicos modernos, la contenerización mediante Docker y la personalización de rutas de aprendizaje mediante Machine Learning, aseguran un producto de alta calidad de software. En consecuencia, EduFinanzas se consolida como una plataforma tecnológica con el potencial de transformar la cultura económica local y promover prácticas de consumo sostenibles mediante la inteligencia de datos y la transparencia del modelo Open Finance. sector financiero, sino también para generar un impacto social positivo y promover prácticas sostenibles dentro de la industria.